\documentclass[11pt]{article}
\usepackage[a4paper,margin=1in]{geometry}
\usepackage{fontspec}
\usepackage[boldfont=false]{xeCJK}
\setCJKmainfont{FandolSong-Regular.otf}
\usepackage{amsmath}
\usepackage{amssymb}
\usepackage{array}
\usepackage{booktabs}
\usepackage{graphicx}
\usepackage{adjustbox}
\usepackage{enumitem}
\setlist[itemize]{topsep=2pt,partopsep=0pt,parsep=0pt,itemsep=2pt,leftmargin=1.4em}
\setlist[enumerate]{topsep=2pt,partopsep=0pt,parsep=0pt,itemsep=2pt,leftmargin=1.6em}
\usepackage{parskip}
\usepackage{fvextra}
\fvset{breaklines=true,breakanywhere=true,fontsize=\small}
\setlength{\parindent}{0pt}

\begin{document}

\begin{center}
  {\LARGE\bfseries Cell membranes and transport}\\[0.35em]
  {\large A Level Biology}
\end{center}
\vspace{0.6em}

\section*{The cell surface membrane}

Every cell is wrapped in a \textbf{cell surface membrane} 细胞膜. We describe its structure with the \textbf{fluid mosaic model} 流动镶嵌模型.

\subsection*{The phospholipid bilayer}

The membrane is built mainly from \textbf{phospholipids} 磷脂. Each phospholipid has a \textbf{hydrophilic} 亲水 ("water-loving") head and two \textbf{hydrophobic} 疏水 ("water-fearing") tails. There is water on both sides of the membrane, so the phospholipids line up in two layers — a \textbf{bilayer} 双层 — with the heads facing the water outside and inside, and the tails hidden in the middle, away from water. This arrangement forms by itself because of those hydrophilic and hydrophobic interactions.

The model is called "fluid" because the phospholipids are not fixed: they slide past each other, so the membrane can move and bend. It is called a "mosaic" because many \textbf{proteins} 蛋白质 are dotted through it, like tiles in a picture.

\subsection*{What floats in the membrane}

\begin{center}
\adjustbox{max width=\linewidth}{%
\begin{tabular}{lll}
\toprule
\textbf{Part} & \textbf{Where it sits} & \textbf{Main roles} \\
\midrule
phospholipids & the two layers & form the basic barrier \\
proteins & through the membrane or on its surface & transport, support and signalling \\
\textbf{carrier proteins} 载体蛋白 & span the membrane & carry specific molecules across \\
\textbf{channel proteins} 通道蛋白 & span the membrane & form water-filled pores for ions to pass \\
\textbf{cholesterol} 胆固醇 & between the phospholipid tails & controls \textbf{fluidity} 流动性 and adds strength \\
\textbf{glycolipids} 糖脂 and \textbf{glycoproteins} 糖蛋白 & carbohydrate chains on the outer surface & cell recognition; some act as \textbf{antigens} 抗原 \\
\bottomrule
\end{tabular}}
\end{center}

So the membrane molecules together give the membrane its \textbf{stability} 稳定性, its fluidity, its \textbf{permeability} 通透性 (control over what gets through), its transport jobs, its signalling jobs, and its cell recognition.

The membrane is \textbf{partially permeable} 半透膜: it lets some substances through easily but blocks others.

\section*{Cell signalling}

Cells talk to each other by \textbf{cell signalling} 细胞信号传递. The main stages are:

\begin{enumerate}
\item a cell \textbf{secretes} a signal chemical called a \textbf{ligand} 配体 (for example a \textbf{hormone} 激素).

\item the ligand is carried (often in the blood) to a \textbf{target cell} 靶细胞.

\item the ligand binds to a specific \textbf{receptor} 受体 on the target cell's surface membrane. The shape of the receptor matches that ligand only. Binding then triggers a particular response inside the target cell.

\end{enumerate}

\section*{Moving substances across the membrane}

There are six processes. Some are \textbf{passive} 被动 (they need no energy 能量), and some are \textbf{active} (they use energy from ATP).

\subsection*{Simple diffusion}

\textbf{Diffusion} 扩散 is the net movement of particles from where they are at a high \textbf{concentration} 浓度 to where they are at a low concentration, until they are spread evenly. \textbf{Simple diffusion} 简单扩散 is when particles pass straight through the bilayer, down the \textbf{concentration gradient} 浓度梯度. Only small or non-polar molecules can do this — such as \textbf{oxygen} 氧气 and \textbf{carbon dioxide} 二氧化碳. It is passive.

\subsection*{Facilitated diffusion}

Charged \textbf{ions} 离子 and large polar molecules (such as \textbf{glucose} 葡萄糖) cannot cross the oily bilayer by themselves. In \textbf{facilitated diffusion} 易化扩散 they cross through channel proteins or carrier proteins, still moving down the concentration gradient. It is also passive.

\subsection*{Osmosis}

\textbf{Osmosis} 渗透 is the diffusion of water across a partially permeable membrane, from a higher \textbf{water potential} 水势 to a lower water potential. It is passive.

\subsection*{Active transport}

\textbf{Active transport} 主动运输 moves a substance \textbf{against} its concentration gradient — from low to high concentration. This needs carrier proteins and energy from ATP.

\subsection*{Endocytosis and exocytosis}

These move large amounts of material in bulk, using ATP.

\begin{itemize}
\item in \textbf{endocytosis} 胞吞作用, the membrane folds inwards around material and pinches off a \textbf{vesicle} 囊泡 to bring it into the cell.

\item in \textbf{exocytosis} 胞吐作用, a vesicle fuses with the membrane and releases its contents outside the cell.

\end{itemize}

\section*{Surface area to volume ratio}

A cell takes in and removes substances across its surface. As an object gets bigger, its \textbf{volume} 体积 grows faster than its \textbf{surface area} 表面积. So the surface area to volume ratio gets \textbf{smaller} as size increases.

For a cube of side $L$:

\[
\text{surface area} = 6L^2, \qquad \text{volume} = L^3, \qquad \text{ratio} = \frac{6}{L}.
\]

A large $L$ gives a small ratio. This is why small cells (and thin, flat shapes) exchange materials quickly, while large cells cannot rely on diffusion alone.

You can show this with \textbf{agar} 琼脂 blocks of different sizes soaked in dye or acid: the smallest block, with the largest surface area to volume ratio, changes colour all the way through fastest. Diffusion across non-living materials can also be studied with \textbf{dialysis tubing} 透析袋 (Visking tubing).

\section*{Water potential and living cells}

\textbf{Water potential} measures how likely water is to leave a solution. Pure water has the highest water potential. Adding a \textbf{solute} 溶质 (a dissolved substance) lowers it. Water always moves by osmosis from a higher to a lower water potential.

To estimate the water potential of plant tissue, you place pieces in sucrose solutions of different water potentials. The solution that causes \textbf{no change} in mass or length has about the same water potential as the tissue.

\subsection*{Effect on plant cells}

\begin{itemize}
\item in a solution of \textbf{higher} water potential (for example \textbf{distilled water} 蒸馏水), water enters the cell. The cell swells and becomes \textbf{turgid} 膨胀, but the strong cell wall stops it bursting.

\item in a solution of \textbf{lower} water potential, water leaves. The cell contents shrink and the membrane pulls away from the cell wall — this is \textbf{plasmolysis} 质壁分离.

\end{itemize}

\subsection*{Effect on animal cells}

Animal cells have no cell wall to protect them.

\begin{itemize}
\item in a solution of higher water potential, water enters and the cell may burst. In a red blood cell this bursting is called \textbf{haemolysis} 溶血.

\item in a solution of lower water potential, water leaves and the cell shrinks.

\end{itemize}


\end{document}
